\documentclass[10pt]{article}
\usepackage[margin=20mm]{geometry}
\usepackage{amsmath,amssymb,mathtools,fontspec,unicode-math,xurl,hyperref}
\IfFontExistsTF{DejaVu Serif}{\setmainfont{DejaVu Serif}}{\setmainfont{Latin Modern Roman}}
\setmathfont{latinmodern-math.otf}
\setlength{\parindent}{0pt}\setlength{\parskip}{6pt}\sloppy
\title{Local Regularity and Finite Time Singularity for the Generalized SQG Equation on the Half-Plane:
Review of the Acquired Sources}
\author{AI-assisted technical review}
\date{}
\begin{document}
\maketitle
Abstract: The manuscript combines a boundary-adapted local theory, two distinct ill-posedness constructions, and a moving-region argument for finite-time singularity. One auxiliary region in the central normal-velocity estimate is defined with a minimum where a maximum is required. Its printed tail integral is infinite, and the claimed positive lower bound is therefore false as written. This is a Major issue in that derivation, with an explicit repair that recovers the subsequent domains. Three Minor corrections concern an overbroad Euler endpoint, a missing critical logarithm, and an ill-posedness prefactor. None is a counterexample to the main theorems. The positive endpoint integral estimates and the repaired proof dependencies were checked. The review distinguishes the pinned 2023 preprint, the separately acquired unversioned author snapshot, and the unavailable journal version; adjudication remains pending.\par
\section*{Sources, quantifiers and proof architecture}
The work is Andrej Zlatoš, Local regularity and finite time singularity for the generalized SQG equation on the half-plane [primary], associated with DOI https://doi.org/10.1215/00127094-2025-0024. The primary actually reviewed is arXiv:2305.02427v2, 5 October 2023, 30 pages. Its matching TeX was read as source, not compiled or executed. A distinct 31-page author-hosted snapshot [alternate] was also acquired and compared. Its PDF metadata says 17 November 2025; that is not an authenticated revision date or a proof of journal equivalence. The publisher version, described by the acquired metadata as Duke Mathematical Journal 174(17), 3493–3533, 15 November 2025, was unavailable behind an access interstitial. The review does not transfer claims to unacquired journal bytes.\par
The 20 manifest-bound files total 2,415,863 bytes and had no intake size or hash discrepancy. The complete primary text and the corresponding scientific sections of the separate author snapshot were covered. No separate supplement or computational repository was located in the bounded searches; the publisher supplement inventory was inaccessible. This is an access limitation, not a mathematical objection.\par
Set \[X_\beta=L^1(\mathbb R\times\mathbb R^+)\cap W^{1,\infty}_\beta,\qquad
 \|\theta\|_{\dot W^{1,\infty}_\beta}
 =\|\partial_1\theta\|_\infty+
   \|\min\{x_2^\beta,1\}\partial_2\theta\|_\infty.\] Theorem 1.1, primary page 4, assumes \(0<\alpha\le1/4\) and \(2\alpha\le\beta\le1-2\alpha\). It states existence, uniqueness, a norm-dependent lifespan and continuation criterion, followed by finite-time singularity for suitable smooth initial data. The definition of classical permits locally Lipschitz solutions satisfying the equation almost everywhere. Smooth data may be nonzero on the boundary; the theorem does not assume compact support strictly inside the open half-plane.\par
Theorem 1.2 must be read together with the opening of Section 3, page 9. For \(\beta<2\alpha\), the construction rules out a solution bounded in \(X_\beta\) on any positive time interval. For \(\beta>1-2\alpha\), it rules out the specified strong \(X_\beta\) continuous dependence for a convergent sequence of data. The latter is not a nonexistence assertion for each datum. The \(H^3\) or \(H^2\) patch-norm blow-up part of Theorem 1.3 is conditional on the corresponding local theory; the manuscript does not prove that entire patch local theory.\par
Section 2 uses the vertical coordinate map \(\Lambda_\beta(x)=(x_1,\lambda_\beta(x_2))\). The identities \(\lambda_\beta'=\kappa_\beta(\lambda_\beta)\), \(\kappa_\beta(s)=\min\{s^\beta,1\}\), convert the weighted norm into an ordinary Lipschitz norm. Section 3 uses boundary shearing at geometrically separated heights. Section 4 propagates a plateau \(\theta=1\) into a growing trapezoid, forcing its distance to the symmetry axis to vanish in finite time.\par
\section*{The printed bad region makes the key integral diverge}
Problem 1 (Primary pages 24–25, definition of \(B_*\) and (4.18)).  With \(a=x_*=x_1=x_2>0\), the manuscript defines \[B_*=\{(p,q):p>a,\ \min\{p,2a\}<q<p+2a\}.
 \tag{A}\] The asserted \(O(a)\) tail of its integral is false. Indeed, the left side of the normal-velocity lower bound (4.18) is \(-\infty\), even with symmetric truncation. The same minimum and tail assertion occur on page 25 of the acquired author snapshot.\par
\subsection*{A direct admissible calculation}
Fix \(0<a\le1/8\) and \(0<\alpha\le1/4\); \(a=1/100,\alpha=1/4\) is one concrete choice. For every \(p\ge1\), the interval \(p/2<q<p\) lies in (A). On this wedge, \[p-a\ge\tfrac78p,\qquad
 (p-a)^2+(q-a)^2\le2p^2.\] Consequently \[\int_{B_*\cap\{p\ge1\}}
       \frac{|p-a|}{[(p-a)^2+(q-a)^2]^{1+\alpha}}\,dp\,dq
 \ge \frac7{16}2^{-1-\alpha}
             \int_1^\infty p^{-2\alpha}\,dp
 =+\infty.
 \tag{B}\] This tests precisely the tail that the proof calls \(O(a)\). It uses admissible parameters and the actual printed domain.\par
In (4.18), the \(B_*\) integral enters with a negative sign. The reflected region \(\widetilde B_*\) enters in the first integral, but its numerator \(p-a\) is negative there. Thus the two divergent contributions have the same adverse sign. The good regions have bounded vertical width and finite integrals for \(\alpha>0\); the bounded rectangle is also integrable because \(\alpha<1/2\). There is no cancellation of opposite infinities. The displayed expression cannot be bounded below by \(Ca\) as claimed.\par
\subsection*{Impact and a complete local repair}
This is Major for the supplied finite-time-singularity derivation. Equation (4.18) supplies \(u_2\ge0\) on the slanted boundary of the moving trapezoid. Together with the negative \(u_1\) estimate, this excludes the first contact and closes the main blow-up argument. The failure is in an essential integral estimate, even though its repair is local.\par
There is also strong counterevidence to a claim that the final theorem is false. Replace (A) by \[B_*^{\rm corr}
   =\{(p,q):p>a,\ \max\{p,2a\}<q<p+2a\}.
 \tag{C}\] The minimum in the separate good region \[G_*=\{(p,q):p>a,\ 0<q<\min\{p,2a\}\}\] must be retained.\par
For \(a<p<1\), reflecting the lower part of the original triangle across \(q=a\) gives the interval \((2a,p+2a)\). If \(p\ge2a\), this splits into \((2a,p)\), already in the upper part of the triangle, and \((p,p+2a)\), which is (C). If \(a<p<2a\), the whole interval is (C). These are exactly the pieces needed for the stated reflection cancellation, with no extra wedge.\par
The corrected strip has height at most \(2a\). Its tail is therefore bounded by \[C_\alpha a\int_1^\infty p^{-1-2\alpha}\,dp
 =O_\alpha(a).\] The subsequent rectangular cancellation adds only \[T=\{(p,q):a<p<2a,\ p<q<2a\}.\] The disjoint union \(B_*^{\rm corr}\cup T\) is \(\{p>a,\ p<q<p+2a\}\). Under \(z=((p-a)/a,(q-a)/a)\), it becomes precisely the printed domain \(B_0=\{z_1>0,z_1<z_2<z_1+2\}\). The reflected corrected strip becomes \[B_0^-=
 \{z_1<-2,\ \max\{-z_1-2,1\}<z_2<-z_1\},\] also precisely the domain in Lemma 4.3. Thus the later, correct maximum provides internal evidence for the repair, and the already stated final positivity lemma applies to it. No change of \(\alpha\), initial data or moving trapezoid is required. This review does not present (B) as a counterexample to the actual velocity inequality or to Theorem 1.1(ii).\par
\section*{Three bounded local corrections}
Problem 2 (The Euler endpoint in Lemma 2.2).  Both acquired texts, page 6, include \(\alpha=0\) in Lemma 2.2. Its estimate (2.3) then asserts \(\partial_1u_2=O(x_2)\) for all permitted data. That assertion is false.\par
Let \(\rho\) be an even nonnegative smooth cutoff, equal to one on \([-1,1]\) and supported in \((-2,2)\). Let \(\sigma\) be nonnegative and smooth on the closed half-line, equal to one on \([0,1]\) and supported in \([0,2)\). Put \[\theta(y_1,y_2)=|y_1|\rho(y_1)\sigma(y_2).\] This fixed function belongs to \(X_0\), and indeed to every \(X_\beta\) under consideration. Its cusp is allowed by the stated Lipschitz hypothesis.\par
Using the source’s kernel normalization and integrating the horizontal derivative through \(\theta\), at \(x=(0,r)\), \[\partial_1u_2(0,r)
 =\int_{y_2>0}y_1\partial_1\theta(y)
 \left\{\frac1{y_1^2+(r-y_2)^2}
        -\frac1{y_1^2+(r+y_2)^2}\right\}\,dy.
 \tag{D}\] On \([-1,1]\times[0,1]\), \(y_1\partial_1\theta=|y_1|\). The remaining contribution is \(O(r)\). Integrating in \(y_1\) shows that the leading part of (D) is \[\int_0^1\left[
 \log\frac{1+(t-r)^2}{1+(t+r)^2}
       +2\log\frac{t+r}{|t-r|}\right]dt
 =4r\log(1/r)+O(r).\] Thus \(r^{-1}|\partial_1u_2(0,r)|\to\infty\). The derivative is continuous for each \(r>0\), so this also contradicts the essential-supremum bound, not merely a value on a line of measure zero. An independent wedge estimate in the evidence package gives the lower bound \((128/841)r\log(1/(4r))-Cr\).\par
This is Minor with respect to the paper’s main claims because every main theorem requires \(\alpha>0\). The proof of (2.3) uses a radial integral that produces \(r^{1-2\alpha}\) for \(\alpha>0\), but \(r\log(1/r)\) for \(\alpha=0\). Restrict this estimate to positive \(\alpha\), or insert the logarithmic Euler boundary modulus. The separate estimate \(u_2=O(x_2)\), derived from \(\partial_2u_2\) and the zero normal boundary value, is not contradicted by this example.\par
Problem 3 (A missing logarithm immediately after (2.8)).  On page 8 of both texts, the proof bounds \[J_{\alpha,\beta}(r)
 =\int_{-2}^{2}|r-y|^{-\beta}|y|^{-2\alpha}\,dy
 \le C_{\alpha,\beta}\max\{r^{1-2\alpha-\beta},1\}.
 \tag{E}\] This fails when \(2\alpha+\beta=1\), including the main endpoint \(\alpha=1/4,\beta=1/2\).\par
For \(0<r<1/4\), at this endpoint the integral over \(2r<y<1\) alone is at least \(\int_{2r}^1dy/y=\log(1/(2r))\), whereas the right-hand maximum in (E) equals one. More precisely, \[J_{1/4,1/2}(r)
 =2\operatorname{arsinh}\sqrt{2/r}+\pi
  +2\operatorname{arcosh}\sqrt{2/r}
 =2\log(1/r)+O(1).\] The needed correction is a logarithmic middle case: bounded below the threshold, \(O(1+|\log r|)\) at it, and \(O(r^{1-2\alpha-\beta})\) above it.\par
The consequence needed for the theorem nevertheless survives. For fixed \(\alpha>0\), multiplication by \(r^{2\alpha}\) controls the critical logarithm; above the threshold it leaves \(r^{1-\beta}\), which is bounded since \(\beta<1\). Hence \(J_{\alpha,\beta}(r)\le C_{\alpha,\beta}r^{-2\alpha}\) for \(0<r\le1\), sufficient for (2.5). The stretched-velocity estimate remains valid when \(\beta\ge2\alpha\), including the closed main endpoint. This is a Minor correction to an intermediate estimate, with an explicit replacement argument.\par
Problem 4 (The prefactor after (3.17), page 13).  The last displayed equality in the first ill-posedness case uses an incorrect positive constant.\par
With the definitions on that page, \(b_n=2^{-2(2\alpha+\beta)n/(1-\beta)}a_n\), and (1.5), direct differentiation gives \[\frac{a_n\lambda_\beta'(2^{-4n-1})}{3b_n}
 =\frac13\left(\frac{1-\beta}{2}\right)^{\beta/(1-\beta)}
           2^{2(2\alpha-\beta)n/(1-\beta)}.
 \tag{F}\] The printed constant instead contains \((2-2\beta)^{\beta/(1-\beta)}\). For the admissible choice \(\alpha=\beta=1/4\), the printed value is larger than (F) by \(2^{2/3}\), for every \(n\). This is an exact counterexample to the displayed equality.\par
Since \(2\alpha-\beta>0\) in this case and the corrected constant is strictly positive, the lower bound still tends to infinity. Neither the crossing-time contradiction nor Theorem 1.2 is refuted. This is Minor; replace the prefactor while retaining the argument. The corresponding text of the author snapshot retains the same factor.\par
Two further algebraic slips can be corrected without changing their subsequent calculations. In the proof of Lemma 4.1 on primary page 21, the parenthetical \(-2y_2+2<0\) has the wrong sign for \(0<y_2<1\); the derivative estimate actually uses \(2y_2-2<0\). On primary page 29, remove the extra outer square in \(x_2^2+(x_2+2)^2=2((x_2+1)^2+1)^2\). At \(x_2=0\) the two printed sides are 4 and 8, whereas the correct right side is \(2((x_2+1)^2+1)\), as used in the next integral. These are local algebraic corrections, not additional failures of the positive integral lemmas.\par
\section*{Checks supporting the intended conclusions}
\subsection*{Local theory and nonlinear uniqueness}
The transformed velocity has bounded derivative when \(2\alpha\le\beta\le1-2\alpha\): the two off-diagonal terms require respectively \(\beta\ge2\alpha\) and \(\beta\le1-2\alpha\). The diagonal terms use the tangential derivative bound and \(u_2=O(x_2)\). Preservation of physical area gives constant \(L^1\) and \(L^\infty\) scalar norms and the Riccati-type gradient estimate (2.9). The logarithmic repair above preserves this reasoning.\par
The compressed uniqueness sentence on page 9 should not be read as inferring nonlinear uniqueness from one fixed Lipschitz velocity. A valid flow comparison is available. Compare two stretched flows from the same initial labels and let their uniform separation be \(\delta(t)\). Their physical coupling preserves area, and \(\lambda_\beta\) is 1-Lipschitz. Split the velocity kernel into a fixed near part and a far part. The near tangential kernel has bounded \(L^1\) norm; the near normal mirror kernel has \(L^1\) norm \(O_\alpha(x_2^{1-2\alpha})\), which remains bounded after division by \(\kappa_\beta(x_2)\) in the allowed range. The far part is estimated after the area-preserving coupling change of variables, using the kernel gradient and the transported \(L^1\) norm. Together these give \[\|\widetilde u^{(1)}-\widetilde u^{(2)}\|_\infty
 \le C\bigl(\operatorname{Lip}\widetilde\theta^{(2)}
                   +\|\theta_0\|_1\bigr)\delta(t),\] and hence \(\delta'\le C(t)\delta\), \(\delta(0)=0\). The full cutoff and boundary estimates are saved separately. Thus no uniqueness defect is inferred from the short exposition. This check does not silently add a strong \(X_\beta\) continuous-dependence theorem to the explicit statement.\par
\subsection*{The two ill-posedness mechanisms}
For the first case, the matching TeX confirms \[a_n=a\,2^{-8n/(1-2\alpha)^2};\] the denominator is squared. The decay of this sequence, the separated supports and the measure-preserving flow control the perturbative velocity terms. The crossing time is bounded by a constant times \(2^{-2(2\alpha-\beta)n/(1-\beta)}\to0\), while the transverse separation cannot exit its comparison interval that quickly. The corrected (F) still contradicts a uniform Lipschitz bound.\par
For the second case, the added \(1/n\) factor makes the truncated data converge in \(X_\beta\). Under the assumed convergence of their solutions, the background-cone and vertical-separation estimates hold uniformly in the truncation index. The crossing times are \[T_n=O\!\left(n^2
       2^{-4n(\beta+2\alpha-1)/(1-\beta)}\right)\to0.\] The concluding Lipschitz lower bound is \(n/3\), as confirmed from the actual page-18 image, not \(n^3\). Its divergence yields the stated failure of continuous dependence. Both exponents have the correct sign on their respective strict parameter ranges.\par
\subsection*{The positive blow-up margins}
The plateau construction is compatible with smooth initial data and odd symmetry. The moving boundary satisfies \(X'=-(50\alpha)^{-1}X^{1-2\alpha}\); the established tangential velocity bound uses \(-(45\alpha)^{-1}X^{1-2\alpha}\), leaving a strict first-contact margin. After the correction (C), the normal estimate is reduced to exactly the printed positive integral lemma.\par
The most delicate tangential integral is also supported by an analytic endpoint certificate. At \(\alpha=1/4\), let \(a_0=2^{1/4}\) and \[R=\frac{4+2^{3/4}-2^{3/2}}{4-2^{-1/4}}
   =\frac{4a_0+2-2a_0^3}{4a_0-1}.\] The elementary chord estimate on page 28 gives \[f(1)=\int_0^1(1+q^2)^{-1/4}\,dq\ge\frac{23}{25}.\] Exact rational fourth-power bounds \[\frac{1189207}{10^6}<a_0<\frac{1189208}{10^6}\] show \(R<271/300\): the equivalent polynomial \(116a_0+871-600a_0^3\) is negative at the lower rational endpoint and decreasing thereafter. Thus \(f(1)>R+1/60\), which yields the required \(I(1)>1/20\). The source’s monotonicity and its separate lower bound for \(I(1)+2g(1)\) then give the uniform \(1/(40\alpha)\) estimate. This check uses exact arithmetic and analytic inequalities, not a plot or unverified numerical quadrature.\par
The remaining normal integral simplifies to \[(2-2^{1-2\alpha})f(1)
 +2^{2-2\alpha}f(1/2)
 -\frac{2^{1-2\alpha}(1-2^{-\alpha})}{1-2\alpha}>0\] in the stipulated range. Its positive lower margin survives the noted algebraic typesetting correction.\par
\section*{Version comparison, assessment and limits}
The acquired author snapshot adds separate Lemmas 4.1 and 4.4 for the pointwise kernel majorant and minorant, and reorganizes the reflection argument. Primary Lemmas 4.1, 4.2 and 4.3 correspond to its Lemmas 4.2, 4.3 and 4.5. This improves the exposition of the same mechanism; the comparison found no changed main parameter range. The critical minimum in \(B_*\), the two local endpoint issues and the prefactor remain in the acquired texts. This is a statement about those files, not the inaccessible journal version.\par
The assessment is one Major proof-step correction and three Minor local corrections. The severity of the first rests on its role in the central finite-time-singularity argument; the complete local repair is separately exhibited. No final theorem has been disproved. The strict ill-posedness ranges, positive integral margins and the weaker velocity consequences of the corrected bounds provide substantive supporting evidence.\par
The review uses the imported whole-plane kernel estimate, patch regularity results and continuation references in their stated roles; it does not claim to have independently reproved every cited theorem or the historical priority claims. The primary executor performed the scientific synthesis, checked the bounded calculations, wrote this report, compiled the current TeX with the existing approved compiler, and actually inspected every current report page. The native helpers `/root/m0191\_rough\_gauge`, `/root/m0191\_pinching` and `/root/m0191\_rates` checked respectively the local theory, ill-posedness constructions and blow-up integrals, with corresponding source-version comparisons. These are bounded AI roles in the same conversation, not human reviews or independent acceptance.\par
Exact source, TeX, compiler, preserved dependencies, logs, PDF and image hashes are bound in the evidence package. Raw render-only and failed-quality observations are preserved separately from actual usable inspections. UI observation and runtime self-description remain distinct; no independently verified backend identity is inferred. This is an initial review with `scientific\_decision=pending`, `human\_reviews=[]`, and no independent acceptance.\par
9 A. Zlatoš, Local regularity and finite time singularity for the generalized SQG equation on the half-plane, arXiv:2305.02427v2, 5 October 2023, 30 pages, https://arxiv.org/pdf/2305.02427v2. PDF SHA-256: `b23939c874c5ce7f29018cd940a7ecd753e0f1a30881dee4f13c732b78a38736`. A. Zlatoš, same title, unversioned author-hosted snapshot, retrieved 26 October 2026, 31 pages, PDF metadata 17 November 2025, https://mathweb.ucsd.edu/\textasciitilde{}zlatos/Papers/sqg\_blowup.pdf. PDF SHA-256: `499d90a9ab04ca72acd220502178045916b0e6c5a602835f885cdccdafccbfbe`. Associated journal DOI: https://doi.org/10.1215/00127094-2025-0024.\par
\end{document}
